\section{USAJMO 2011}
        \begin{enumerate}
        	\item (\textit{USAJMO 2011}) Find, with proof, all positive integers $n$ for which $2^n + 12^n + 2011^n$ is a perfect square.


 



\textbf{Solution 1}

Let $2^n + 12^n + 2011^n = x^2$. Then $(-1)^n + 1 \equiv x^2 \pmod {3}$. Since all perfect squares are congruent to 0 or 1 modulo 3, this means that n must be odd. Proof by Contradiction:We wish to show that the only value of $n$ that satisfies is $n = 1$. Assume that $n \ge 2$. Then consider the equation$2^n + 12^n = x^2 - 2011^n$. From modulo 2, we easily know x is odd. Let $x = 2a + 1$, where a is an integer. $2^n + 12^n = 4a^2 + 4a + 1 - 2011^n$. Dividing by 4,$2^{n-2} + 3^n \cdot 4^{n-1} = a^2 + a + \dfrac {1}{4} (1 - 2011^n)$. Since $n \ge 2$, $n-2 \ge 0$, so $2^{n-2}$ similarly, the entire LHS is an integer, and so are $a^2$ and $a$. Thus, $\dfrac {1}{4} (1 - 2011^n)$ must be an integer. Let $\dfrac {1}{4} (1 - 2011^n) = k$. Then we have $1- 2011^n = 4k$. $1- 2011^n \equiv 0 \pmod {4}$. $(-1)^n \equiv 1 \pmod {4}$. Thus, n is even. However, it has already been shown that $n$ must be odd.  This is a contradiction.  Therefore, $n$ is not greater than or equal to 2, and must hence be less than 2.  The only positive integer less than 2 is 1.

 



\textbf{Solution 2}

If $n = 1$, then $2^n + 12^n + 2011^n = 2025 = 45^2$, a perfect square. 

 If $n > 1$ is odd, then $2^n + 12^n + 2011^n \equiv 0 + 0 + (-1)^n \equiv 3 \pmod{4}$. 

 Since all perfect squares are congruent to $0,1 \pmod{4}$, we have that $2^n+12^n+2011^n$ is not a perfect square for odd $n > 1$. 

 If $n = 2k$ is even, then $(2011^k)^2 < 2^{2k}+12^{2k}+2011^{2k}$ $= 4^k + 144^k + 2011^{2k} <$ $2011^k + 2011^k + 2011^{2k} < (2011^k+1)^2$. 

 Since $(2011^k)^2 < 2^n+12^n+2011^n < (2011^k+1)^2$, we have that $2^n+12^n+2011^n$ is not a perfect square for even $n$. 

 Thus, $n = 1$ is the only positive integer for which $2^n + 12^n + 2011^n$ is a perfect square.

 



\textbf{Solution 3}

Looking at residues mod 3, we see that $n$ must be odd, since even values of $n$ leads to $2^n + 12^n + 2011^n \equiv 2 \pmod{3}$.  Also as shown in solution 2, for $n>1$, $n$ must be even.  Hence, for $n>1$, $n$ can neither be odd nor even.  The only possible solution is then $n=1$, which indeed works.

 



\textbf{Solution 4}

Take the whole expression mod 12. Note that the perfect squares can only be of the form 0, 1, 4 or 9 (mod 12). Note that since the problem is asking for positive integers, $12^n$ is always divisible by 12, so this will be disregarded in this process. If $n$ is even, then $2^n \equiv 4 \pmod{12}$ and $2011^n \equiv 7^n \equiv 1 \pmod {12}$. Therefore, the sum in the problem is congruent to $5 \pmod {12}$, which cannot be a perfect square. Now we check the case for which $n$ is an odd number greater than 1. Then $2^n \equiv 8 \pmod{12}$ and $2011^n \equiv 7^n \equiv 7 \pmod {12}$. Therefore, this sum would be congruent to $3 \pmod {12}$, which cannot be a perfect square. The only case we have not checked is $n=1$. If $n=1$, then the sum in the problem is equal to $2+12+2011=2025=45^2$. Therefore the only possible value of $n$ such that $2^n+12^n+2011^n$ is a perfect square is $n=1$.

 



\textbf{Solution 5}

We will first take the expression modulo $3$. We get $2^n+12^n+2011^n \equiv -1^n+1^n \pmod 3$. 

 Lemma 1: All perfect squares are equal to $0$ or $1$ modulo $3$.We can prove this by testing the residues modulo $3$. We have $0^2 \equiv 0 \pmod 3$, $1^2 \equiv 1 \pmod 3$, and $2^2 \equiv 1 \pmod 3$, so the lemma is true.

 We know that if $n$ is odd, $-1^n+1^n \equiv 0 \pmod 3$, which satisfies the lemma's conditions. However, if $n$ is even, we get $2 \pmod 3$, which does not satisfy the lemma's conditions. So, we can conclude that $n$ is odd.

 Now, we take the original expression modulo $4$. For right now, we will assume that $n>1$, and test $n=1$ later. For $n>1$, $2^n \equiv 0 \pmod 4$, so $2^n+12^n+2011^n=-1^n \pmod 4$.

 Lemma 2: All perfect squares are equal to $0$ or $1$ modulo $4$.We can prove this by testing the residues modulo $4$. We have $0^2 \equiv 0 \pmod 4$, $1^2 \equiv 1 \pmod 4$, $2^2 \equiv 0 \pmod 4$, and $3^2 \equiv 1 \pmod 4$, so the lemma is true.

 We know that if $n$ is even, $-1^n \equiv 0 \pmod 4$, which satisfies the lemma's conditions. However, if $n$ is odd, $-1^n \equiv -1 \equiv 3 \pmod 4$, which does not satisfy the lemma's conditions. Therefore, $n$ must be even.

 However, a number cannot be even and odd at the same time, so this is impossible. Now, we only have to test $n=1$. We know that $2^1+12^1+2011^1=45^2$, so the only integer is $\boxed{n=1}$.

 



\textbf{Solution 6}

The answer is $n=1$, which is easily verified to be a valid integer $n$.Notice that 
 \[2^n+12^n+2011^n\equiv 2^n+7^n \pmod{12}.\] Then for $n\geq 2$, we have $2^n+7^n\equiv 3,5 \pmod{12}$ depending on the parity of $n$. But perfect squares can only be $0,1,4,9\pmod{12}$, contradiction. Therefore, we are done. $\blacksquare$

 


	\item (\textit{USAJMO 2011}) Let $a$, $b$, $c$ be positive real numbers such that $a^2 + b^2 + c^2 + (a + b + c)^2 \le 4$.  Prove that

 \[\frac{ab + 1}{(a + b)^2} + \frac{bc + 1}{(b + c)^2} + \frac{ca + 1}{(c + a)^2} \ge 3.\]


 



\textbf{Solution 1}

Since
 \begin{align*} (a+b)^2 + (b+c)^2 + (c+a)^2 &= 2(a^2 + b^2 + c^2 + ab + bc + ca) \\ 	&= a^2 + b^2 + c^2 + (a + b + c)^2, \end{align*}
it is natural to consider a change of variables:
 \begin{align*} \alpha &= b + c \\ \beta &= c + a \\ \gamma &= a + b \end{align*}
with the inverse mapping given by:
 \begin{align*} a &= \frac{\beta + \gamma - \alpha}2 \\ b &= \frac{\alpha + \gamma - \beta}2 \\ c &= \frac{\alpha + \beta - \gamma}2 \end{align*}
With this change of variables, the constraint becomes
 \[\alpha^2 + \beta^2 + \gamma^2 \le 4,\]while the left side of the inequality we need to prove is now
 \begin{align*} & \frac{\gamma^2 - (\alpha - \beta)^2 + 4}{4\gamma^2} + \frac{\alpha^2 - (\beta - \gamma)^2 + 4}{4\alpha^2} + \frac{\beta^2 - (\gamma - \alpha)^2 + 4}{4\beta^2} \ge \\ & \frac{\gamma^2 - (\alpha - \beta)^2 + \alpha^2 + \beta^2 + \gamma^2}{4\gamma^2} + \frac{\alpha^2 - (\beta - \gamma)^2 + \alpha^2 + \beta^2 + \gamma^2}{4\alpha^2} + \frac{\beta^2 - (\gamma - \alpha)^2 + \alpha^2 + \beta^2 + \gamma^2}{4\beta^2} = \\ & \frac{2\gamma^2 + 2\alpha\beta}{4\gamma^2} + \frac{2\alpha^2 + 2\beta\gamma}{4\alpha^2} + \frac{2\beta^2 + 2\gamma\alpha}{4\beta^2} = \\ & \frac32 + \frac{\alpha\beta}{2\gamma^2} + \frac{\beta\gamma}{2\alpha^2} + \frac{\gamma\alpha}{2\beta^2}. \end{align*}
Therefore it remains to prove that
 \[\frac{\alpha\beta}{2\gamma^2} + \frac{\beta\gamma}{2\alpha^2} + \frac{\gamma\alpha}{2\beta^2} \ge \frac32.\]We note that the product of the three (positive) terms is 1/8, therefore by AM-GM their mean is at least 1/2, and thus their sum is at least 3/2 and we are done.

 



\textbf{Solution 2}

Rearranging the condition yields that
 \[a^2 + b^2 + c^2 +ab+bc+ac \le 2\]

 Now note that
 \[\frac{2ab+2}{(a+b)^2} \ge \frac{2ab+a^2 + b^2 + c^2 +ab+bc+ac}{(a+b)^2}=\frac{(a+b)^2 + (c+a)(c+b)}{(a+b)^2}\]

 Summing this for all pairs of $\{ a,b,c \}$ gives that
 \[\sum_{cyc} \frac{2ab+2}{(a+b)^2} \ge 3+ \sum_{cyc}\frac{(c+a)(c+b)}{(a+b)^2} \ge 6\]

 By AM-GM. Dividing by $2$ gives the desired inequality.

 



\textbf{Solution 3}

Given $a^2+b^2+c^2+(a+b+c)^2 \leq 4$, prove that
 \[\frac{ab+1}{(a+b)^2}+\frac{bc+1}{(b+c)^2}+\frac{ac+1}{(a+c)^2} \geq 3.\]

 From the given constraint $a^2+b^2+c^2+(a+b+c)^2 \leq 4$, we can expand to obtain:
 \[a^2+b^2+c^2+a^2+b^2+c^2+2ab+2bc+2ca \leq 4\]
 \[2(a^2+b^2+c^2)+2(ab+bc+ca) \leq 4\]
 \[\frac{a^2+b^2+c^2}{2}+\frac{ab+bc+ca}{2} \leq 1.\]

 Now, observe that we can rewrite the left-hand side of our inequality as:\begin{align*}&\sum_{\text{cyc}} \frac{ab+1}{(a+b)^2} \\&= \sum_{\text{cyc}} \frac{ab+\frac{a^2+b^2+c^2}{2}+\frac{ab+bc+ca}{2}}{(a+b)^2} \\&= \sum_{\text{cyc}} \left[\frac{ab+\frac{a^2+b^2}{2}}{(a+b)^2} + \frac{\frac{c^2}{2}+\frac{bc+ca}{2}}{(a+b)^2}\right] \\&= \sum_{\text{cyc}} \left[\frac{ab+\frac{a^2+b^2}{2}}{(a+b)^2}\right] + \sum_{\text{cyc}} \left[\frac{\frac{a^2+c^2}{2}}{(a+c)^2} + \frac{\frac{ac+bc}{2}}{(a+b)^2}\right].\end{align*}


 We now evaluate each part separately.Note that 
 \[\frac{ab+\frac{a^2+b^2}{2}}{(a+b)^2} = \frac{\frac{(a+b)^2}{2}}{(a+b)^2} = \frac{1}{2}.\]Therefore, 
 \[\sum_{\text{cyc}} \frac{ab+\frac{a^2+b^2}{2}}{(a+b)^2} = \frac{3}{2}.\]By the RMS-AM inequality, we have $(a+c)^2 \leq 2(a^2+c^2)$, which gives $\frac{1}{(a+c)^2} \geq \frac{1}{2(a^2+c^2)}$. Thus,\begin{align*}\sum_{\text{cyc}} \frac{\frac{a^2+c^2}{2}}{(a+c)^2} &\geq \sum_{\text{cyc}} \frac{\frac{a^2+c^2}{2}}{2(a^2+c^2)} \\&= \sum_{\text{cyc}} \frac{1}{4} = \frac{3}{4}.\end{align*}
The remaining terms can be written as:\begin{align*}\sum_{\text{cyc}} \frac{\frac{ac+bc}{2}}{(a+b)^2} &= \frac{1}{2}\sum_{\text{cyc}} \frac{c(a+b)}{(a+b)^2} \\&= \frac{1}{2}\sum_{\text{cyc}} \frac{c}{a+b} \\&= \frac{1}{2}\left(\frac{a}{b+c}+\frac{b}{a+c}+\frac{c}{a+b}\right).\end{align*}
By Nesbitt's inequality, 
 \[\frac{a}{b+c}+\frac{b}{a+c}+\frac{c}{a+b} \geq \frac{3}{2},\]so this part contributes at least $\frac{3}{4}$.Adding all three parts together:
 \[\frac{ab+1}{(a+b)^2}+\frac{bc+1}{(b+c)^2}+\frac{ac+1}{(a+c)^2} \geq \frac{3}{2}+\frac{3}{4}+\frac{3}{4} = 3.\]

 Therefore, the inequality is proven. $\square$

 


	\item (\textit{USAJMO 2011}) For a point $P = (a,a^2)$ in the coordinate plane, let $\ell(P)$ denote the line passing through $P$ with slope $2a$.  Consider the set of triangles with vertices of the form $P_1 = (a_1, a_1^2)$, $P_2 = (a_2, a_2^2)$, $P_3 = (a_3, a_3^2)$, such that the intersections of the lines $\ell(P_1)$, $\ell(P_2)$, $\ell(P_3)$ form an equilateral triangle $\Delta$.  Find the locus of the center of $\Delta$ as $P_1 P_2 P_3$ ranges over all such triangles.


 



\textbf{Solution 1}

Note that the lines $l(P_1), l(P_2), l(P_3)$ are 
 \[y=2a_1x-a_1^2, y=2a_2x-a_2^2, y=2a_3x-a_3^2,\] respectively. It is easy to deduce that the three points of intersection are 
 \[\left(\frac{a_1+a_2}{2},a_1a_2\right),\left(\frac{a_2+a_3}{2},a_2a_3\right), \left(\frac{a_3+a_1}{2},a_3a_1\right).\] The slopes of each side of this equilateral triangle are 
 \[2a_1,2a_2,2a_3,\] and we want to find the locus of 
 \[\left(\frac{a_1+a_2+a_3}{3},\frac{a_1a_2+a_2a_3+a_3a_1}{3}\right).\] Define the three complex numbers $w_n = 1+2a_ni$ for $n=1,2,3$. Then note that the slope - that is, the imaginary part divided by the real part - of all $w_n^3$ is constant, say it is $k$. Then for $n=1,2,3$,

 
 \begin{align*} \frac{\Im(w_n^3)}{\Re(w_n^3)} &= \frac{\Im((1+2a_ni)^3)}{\Re((1+2a_ni)^3)}\\ &= \frac{\Im(1+6a_ni-12a_n^2-8a_n^3i)}{\Re(1+6a_ni-12a_n^2-8a_n^3i)}\\ &= \frac{6a_n-8a_n^3}{1-12a_n^2}\\ &= k.\\ \end{align*}


 Rearranging, we get that
 \[8a_n^3 -12ka_n^2-6a_n+k=0,\]or
 \[a_n^3-\frac{3ka_n^2}2-\frac{3a_n}4+\frac k8=0.\]Note that this is a cubic, and the roots are $a_1,a_2$ and $a_3$ which are all distinct, and so there are no other roots. Using Vieta's, we get that
 \[a_1+a_2+a_3=\frac{3k}2,\]and
 \[a_1a_2+a_2a_3+a_3a_1=-\frac34.\]Obviously all values of $k$ are possible, and so our answer is the line 
 \[\boxed{y=-\frac{1}{4}.}\] $\blacksquare$

 



\textbf{Solution 2}

Note that all the points $P=(a,a^2)$ belong to the parabola $y=x^2$ which we will denote $p$. This parabola has a focus $F=\left(0,\frac{1}{4}\right)$ and directrix $y=-\frac{1}{4}$ which we will denote $d$. We will prove that the desired locus is $d$.

 First note that for any point $P$ on $p$, the line $\ell(P)$ is the tangent line to $p$ at $P$. This is because $\ell(P)$ contains $P$ and because $[\frac{d}{dx}] x^2=2x$. If you don't like calculus, you can also verify that $\ell(P)$ has equation $y=2a(x-a)+a^2$ and does not intersect $y=x^2$ at any point besides $P$. Now for any point $P$ on $p$ let $P'$ be the foot of the perpendicular from $P$ onto $d$. Then by the definition of parabolas, $PP'=PF$. Let $q$ be the perpendicular bisector of $\overline{P'F}$. Since $PP'=PF$, $q$ passes through $P$. Suppose $K$ is any other point on $q$ and let $K'$ be the foot of the perpendicular from $K$ to $d$. Then in right $\Delta KK'P'$, $KK'$ is a leg and so $KK'<KP'=KF$. Therefore $K$ cannot be on $p$. This implies that $q$ is exactly the tangent line to $p$ at $P$, that is $q=\ell(P)$. So we have proved Lemma 1: If $P$ is a point on $p$ then $\ell(P)$ is the perpendicular bisector of $\overline{P'F}$.

 We need another lemma before we proceed. Lemma 2: If $F$ is on the circumcircle of $\Delta XYZ$ with orthocenter $H$, then the reflections of $F$ across $\overleftrightarrow{XY}$, $\overleftrightarrow{XZ}$, and $\overleftrightarrow{YZ}$ are collinear with $H$.

 Proof of Lemma 2: Say the reflections of $F$ and $H$ across $\overleftrightarrow{YZ}$ are $C'$ and $J$, and the reflections of $F$ and $H$ across $\overleftrightarrow{XY}$ are $A'$ and $I$. Then we angle chase $\angle JYZ=\angle HYZ=\angle HXZ=\angle JXZ=m(JZ)/2$ where $m(JZ)$ is the measure of minor arc $JZ$ on the circumcircle of $\Delta XYZ$. This implies that $J$ is on the circumcircle of $\Delta XYZ$, and similarly $I$ is on the circumcircle of $\Delta XYZ$. Therefore $\angle C'HJ=\angle FJH=m(XF)/2$, and $\angle A'HX=\angle FIX=m(FX)/2$. So $\angle C'HJ = \angle A'HX$. Since $J$, $H$, and $X$ are collinear it follows that $C'$, $H$ and $A'$ are collinear. Similarly, the reflection of $F$ over $\overleftrightarrow{XZ}$ also lies on this line, and so the claim is proved.

 Now suppose $A$, $B$, and $C$ are three points of $p$ and let $\ell(A)\cap\ell(B)=X$, $\ell(A)\cap\ell(C)=Y$, and $\ell(B)\cap\ell(C)=Z$. Also let $A''$, $B''$, and $C''$ be the midpoints of $\overline{A'F}$, $\overline{B'F}$, and $\overline{C'F}$ respectively. Then since $\overleftrightarrow{A''B''}\parallel \overline{A'B'}=d$ and $\overleftrightarrow{B''C''}\parallel \overline{B'C'}=d$, it follows that $A''$, $B''$, and $C''$ are collinear. By Lemma 1, we know that $A''$, $B''$, and $C''$ are the feet of the altitudes from $F$ to $\overline{XY}$, $\overline{XZ}$, and $\overline{YZ}$. Therefore by the Simson Line Theorem, $F$ is on the circumcircle of $\Delta XYZ$. If $H$ is the orthocenter of $\Delta XYZ$, then by Lemma 2, it follows that $H$ is on $\overleftrightarrow{A'C'}=d$. It follows that the locus described in the problem is a subset of $d$.

 Since we claim that the locus described in the problem is $d$, we still need to show that for any choice of $H$ on $d$ there exists an equilateral triangle with center $H$ such that the lines containing the sides of the triangle are tangent to $p$. So suppose $H$ is any point on $d$ and let the circle centered at $H$ through $F$ be $O$. Then suppose $A$ is one of the intersections of $d$ with $O$. Let $\angle HFA=3\theta$, and construct the ray through $F$ on the same halfplane of $\overleftrightarrow{HF}$ as $A$ that makes an angle of $2\theta$ with $\overleftrightarrow{HF}$. Say this ray intersects $O$ in a point $B$ besides $F$, and let $q$ be the perpendicular bisector of $\overline{HB}$. Since $\angle HFB=2\theta$ and $\angle HFA=3\theta$, we have $\angle BFA=\theta$. By the inscribed angles theorem, it follows that $\angle AHB=2\theta$. Also since $HF$ and $HB$ are both radii, $\Delta HFB$ is isosceles and $\angle HBF=\angle HFB=2\theta$. Let $P_1'$ be the reflection of $F$ across $q$. Then $2\theta=\angle FBH=\angle C'HB$, and so $\angle C'HB=\angle AHB$. It follows that $P_1'$ is on $\overleftrightarrow{AH}=d$, which means $q$ is the perpendicular bisector of $\overline{FP_1'}$.

 Let $q$ intersect $O$ in points $Y$ and $Z$ and let $X$ be the point diametrically opposite to $B$ on $O$. Also let $\overline{HB}$ intersect $q$ at $M$. Then $HM=HB/2=HZ/2$. Therefore $\Delta HMZ$ is a $30-60-90$ right triangle and so $\angle ZHB=60^{\circ}$. So $\angle ZHY=120^{\circ}$ and by the inscribed angles theorem, $\angle ZXY=60^{\circ}$. Since $ZX=ZY$ it follows that $\Delta ZXY$ is and equilateral triangle with center $H$.

 By Lemma 2, it follows that the reflections of $F$ across $\overleftrightarrow{XY}$ and $\overleftrightarrow{XZ}$, call them $P_2'$ and $P_3'$, lie on $d$. Let the intersection of $\overleftrightarrow{YZ}$ and the perpendicular to $d$ through $P_1'$ be $P_1$, the intersection of $\overleftrightarrow{XY}$ and the perpendicular to $d$ through $P_2'$ be $P_2$, and the intersection of $\overleftrightarrow{XZ}$ and the perpendicular to $d$ through $P_3'$ be $P_3$. Then by the definitions of $P_1'$, $P_2'$, and $P_3'$ it follows that $FP_i=P_iP_i'$ for $i=1,2,3$ and so $P_1$, $P_2$, and $P_3$ are on $p$. By lemma 1, $\ell(P_1)=\overleftrightarrow{YZ}$, $\ell(P_2)=\overleftrightarrow{XY}$, and $\ell(P_3)=\overleftrightarrow{XZ}$. Therefore the intersections of $\ell(P_1)$, $\ell(P_2)$, and $\ell(P_3)$ form an equilateral triangle with center $H$, which finishes the proof.--Killbilledtoucan

 



\textbf{Solution 3}

Note that the lines $l(P_1), l(P_2), l(P_3)$ are 
 \[y=2a_1x-a_1^2, y=2a_2x-a_2^2, y=2a_3x-a_3^2,\] respectively. It is easy to deduce that the three points of intersection are 
 \[\left(\frac{a_1+a_2}{2},a_1a_2\right),\left(\frac{a_2+a_3}{2},a_2a_3\right), \left(\frac{a_3+a_1}{2},a_3a_1\right).\] The slopes of each side of this equilateral triangle are 
 \[2a_1,2a_2,2a_3,\] and we want to find the locus of 
 \[\left(\frac{a_1+a_2+a_3}{3},\frac{a_1a_2+a_2a_3+a_3a_1}{3}\right).\] We know that 
 \[2a_1=\tan(\theta), 2a_2=\tan (\theta + 120), 2a_3=\tan (\theta-120)\] for some $\theta.$ Therefore, we can use the tangent addition formula to deduce 
 \[\frac{a_1+a_2+a_3}{3}=\frac{\tan(\theta)+\tan (\theta + 120)+\tan (\theta-120)}{6}=\frac{3\tan\theta-\tan^3\theta}{2-6\tan^2\theta}\] and 
 \begin{align*} \frac{a_1a_2+a_2a_3+a_3a_1}{3}&=\frac{\tan\theta (\tan(\theta-120)+\tan(\theta+120))+\tan(\theta-120)\tan(\theta+120)}{12}\\ &=\frac{9\tan^2\theta-3}{12(1-3\tan^2\theta)}\\ &=-\frac{1}{4}.\end{align*}
 Now we show that $\frac{a_1+a_2+a_3}{3}$ can be any real number. Let's say 
 \[\frac{3\tan\theta-\tan^3\theta}{2-6\tan^2\theta}=k\] for some real number $k.$ Multiplying both sides by $2-\tan^2\theta$ and rearranging yields a cubic in $\tan\theta.$ Clearly this cubic has at least one real solution. As $\tan \theta$ can take on any real number, all values of $k$ are possible, and our answer is the line 
 \[\boxed{y=-\frac{1}{4}.}\] Of course, as the denominator could equal 0, we must check $\tan \theta=\pm \frac{1}{\sqrt{3}}.$ 
 \[3\tan \theta-\tan^3\theta=k(2-6\tan^2\theta).\] The left side is nonzero, while the right side is zero, so these values of $\theta$ do not contribute to any values of $k.$ So, our answer remains the same. $\blacksquare$ ~ Benq

 \textbf{Work in progress: Solution 4 (Clean algebra)} 
                    \begin{figure}[H]
                        \centering
                        \includegraphics[width=0.5\textwidth]{Solution/USAJMO_I_2011_3_Solution_Figure_1.png}
                    \end{figure}
                    

 It can be easily shown that the center of $\triangle$ has coordinates $\left(\frac{a_{1}+a_{2}+a_{3}}{3},\frac{a_{1}a_{2}+a_{2}a_{3}+a_{3}a_{1}}{3}\right)$.

 Without loss of generality, let $a_{1}>a_{2}>a_{3}$. Notice that $\ell(P_2)$ is a $60^{\circ}$ clockwise rotation of $\ell(P_1)$, $\ell(P_3)$ is a $60^{\circ}$ clockwise rotation of $\ell(P_2)$, and $\ell(P_1)$ is a $60^{\circ}$ clockwise rotation of $\ell(P_3)$. By definition, $\arctan(2a_{i})$ is the (directed) angle from the x-axis to $\ell(P_{i})$. Remember that the range of $\arctan(x)$ is $(-90^{\circ},90^{\circ})$. We have 
 \begin{align*}\arctan(2a_{1})-\arctan(2a_{2})&=60^{\circ}\\\arctan(2a_{2})-\arctan(2a_{3})&=60^{\circ}\\\arctan(2a_{3})-\arctan(2a_{1})&=-120^{\circ}.\end{align*}


 Taking the tangent of both sides of each equation and rearranging, we get 
 \begin{align*}2a_{1}-2a_{2}&=\sqrt{3}(1+4a_{1}a_{2})\\2a_{2}-2a_{3}&=\sqrt{3}(1+4a_{2}a_{3})\\2a_{3}-2a_{1}&=\sqrt{3}(1+4a_{3}a_{1}).\end{align*}
 We add these equations to get 
 \[\sqrt{3}(3+4(a_{1}a_{2}+a_{2}a_{3}+a_{3}a_{1}))=0.\] We solve for $a_{1}a_{2}+a_{2}a_{3}+a_{3}a_{1}$ to get 
 \[a_{1}a_{2}+a_{2}a_{3}+a_{3}a_{1}=-\frac{3}{4}.\] So, the y-coordinate of $\triangle$ is $-\frac{1}{4}$.

 We will prove that the x-coordinate of $\triangle$ can be any real number. If $2a_{1}$ tends to infinity, then $2a_{2}$ tends to $\frac{\sqrt{3}}{2}$ and $2a_{3}$ tends to $-\frac{\sqrt{3}}{2}$. So, $\frac{a_{1}+a_{2}+a_{3}}{3}$ can be arbitrarily large. Similarly, if we let $2a_{3}$ tend to negative infinity, then $2a_{1}$ tends to $\frac{\sqrt{3}}{2}$ and $2a_{2}$ tends to $-\frac{\sqrt{3}}{2}$. So, $\frac{a_{1}+a_{2}+a_{3}}{3}$ can be arbitrarily small. Since $\frac{a_{1}+a_{2}+a_{3}}{3}$ is continuous, it can take any real value. So, the locus is the line $y=-\frac{1}{4}$.

 


	\item (\textit{USAJMO 2011}) A \textit{word} is defined as any finite string of letters.  A word is a \textit{palindrome} if it reads the same backwards as forwards.  Let a sequence of words $W_0$, $W_1$, $W_2$, $\dots$ be defined as follows: $W_0 = a$, $W_1 = b$, and for $n \ge 2$, $W_n$ is the word formed by writing $W_{n - 2}$ followed by $W_{n - 1}$.  Prove that for any $n \ge 1$, the word formed by writing $W_1$, $W_2$, $\dots$, $W_n$ in succession is a palindrome.


 



\textbf{Solution}

Let $r$ be the reflection function on the set of words, namely $r(a_1\dots a_n) = a_n \dots a_1$ for all words $a_1 \dots a_n$, $n\ge 1$. Then the following property is evident (e.g. by mathematical induction):

 $r(w_1 \dots w_k) = r(w_k) \dots r(w_1)$, for any words $w_1, \dots, w_k$, $k \ge 1$.

 We use mathematical induction to prove the statement of the problem. First, $W_1 = b$, $W_1W_2 = bab$, $W_1W_2W_3 = babbab$  are palindromes. Second, suppose $n\ge 3$, and that the words $W_1 W_2 \dots W_k$  ($k = 1$, $2$, $\dots$, $n$) are all palindromes, i.e. $r(W_1W_2\dots W_k) = W_1W_2\dots W_k$. Now, consider the word $W_1 W_2 \dots W_{n+1}$:

 
 \[r(W_1 W_2 \dots W_{n+1}) = r(W_{n+1}) r(W_1 W_2 \dots W_{n-2}W_{n-1}W_n)\]

 
 \[= r(W_{n-1}W_n) W_1 W_2 \dots W_{n-2} W_{n-1} W_n\]

 
 \[= r(W_{n-1}W_n) r(W_1 W_2 \dots W_{n-2}) W_{n+1}\]

 
 \[= r(W_1 W_2 \dots W_{n-2} W_{n-1}W_n) W_{n+1}\]

 
 \[= W_1W_2\dots W_n W_{n+1}.\] 

 By the principle of mathematical induction, the statement of the problem is proved. \href{https://artofproblemsolving.com/wiki/index.php?title=User:Lightest}{Lightest} 21:54, 1 April 2012 (EDT)

 


	\item (\textit{USAJMO 2011}) Points $A$, $B$, $C$, $D$, $E$ lie on a circle $\omega$ and point $P$ lies outside the circle.  The given points are such that (i) lines $PB$ and $PD$ are tangent to $\omega$, (ii) $P$, $A$, $C$ are collinear, and (iii) $\overline{DE} \parallel \overline{AC}$.  Prove that $\overline{BE}$ bisects $\overline{AC}$.


 



\textbf{Solution}

Since $\overline{DE} \parallel \overline{AC}$, we have that $AD=CE\rightarrow \angle ABD=\angle CBE$. But it is well known that $ABCD$ is a harmonic quadrilateral, thus $\overline{BD}$ is a symmedian of triangle $ABC$, from which it follows that $\overline{BE}$ is a median of $\triangle ABC$. $\blacksquare$

 



\textbf{Solution 1}

Connect segment PO, and name the interaction of PO and the circle as point M. 

 Since PB and PD are tangent to the circle, it's easy to see that M is the midpoint of arc BD.

 $\angle$ BOA = 1/2 arc AB + 1/2 arc CE

 Since AC // DE, arc AD = arc CE, 

 thus, $\angle$ BOA = 1/2 arc AB + 1/2 arc AD = 1/2 arc BD = arc BM = $\angle$ BOM

 Therefore, PBFO is cyclic, $\angle$ PFO = $\angle$ OBP = 90°, AF = AC (F is the interaction of BE and AC)

 BE bisects AC, proof completed!

 



\textbf{Solution 2}

Let $O$ be the center of the circle, and let $X$ be the intersection of $AC$ and $BE$. Let $\angle OPA$ be $x$ and $\angle OPD$ be $y$.

 $\angle OPB = \angle OPD = y$, $\angle BED = \frac{\angle DOB}{2} = 90-y$, $\angle ODE = \angle PDE - 90 = 90-x-y$ $\angle OBE = \angle PBE - 90 = x = \angle OPA$

 Thus $PBXO$ is a cyclic quadrilateral and $\angle OXP = \angle OBP = 90$ and so $X$ is the midpoint of chord $AC$.

 



\textbf{Solution 3}

This is the solution from EGMO Problem 1.43 page 242

 Let $O$ be the center of the circle, and let $M$ be the midpoint of $AC$.  Let $\theta$ denote the circle with diameter $OP$.  Since $\angle OBP = \angle OMP = \angle ODP = 90^\circ$, $B$, $D$, and $M$ all lie on $\theta$.

 
                    \begin{figure}[H]
                        \centering
                        \includegraphics[width=0.5\textwidth]{Solution/USAJMO_I_2011_5_Solution_Figure_1.png}
                    \end{figure}
                    

 Since quadrilateral $BOMP$ is cyclic, $\angle BMP = \angle BOP$.  Triangles $BOP$ and $DOP$ are congruent, so $\angle BOP = \angle BOD/2 = \angle BED$, so $\angle BMP = \angle BED$. Because $AC$ and $DE$ are parallel, $M$ lies on $BE$ (using Euclid's Parallel Postulate).

 -Evan Chen (vEnhance)

 



\textbf{Solution 4}

Note that by Lemma 9.9 of EGMO, $(A,C;B,D)$ is a harmonic bundle. We project through $E$ onto $\overline{AC}$,
 \[-1=(A,C;B,D)\stackrel{E}{=}(A,C;M,P_{\infty})\]Where $P_{\infty}$ is the point at infinity for parallel lines $\overline{DE}$ and $\overline{AC}$. Thus, we get $\frac{MA}{MC}=-1$, and $M$ is the midpoint of $AC$. ~novus677

 \textbf{The power of angle chasing=} 


	\item (\textit{USAJMO 2011}) Consider the assertion that for each positive integer $n \ge 2$, the remainder upon dividing $2^{2^n}$ by $2^n - 1$ is a power of 4.  Either prove the assertion or find (with proof) a counterexample.


 



\textbf{Solution}

We will show that $n = 25$ is a counter-example.

 Since $\textstyle 2^n \equiv 1 \pmod{2^n - 1}$, we see that for any integer $k$, $\textstyle 2^{2^n} \equiv 2^{(2^n - kn)} \pmod{2^n-1}$. Let $0 \le m < n$ be the residue of $2^n \pmod n$. Note that since $\textstyle m < n$ and $\textstyle n \ge 2$, necessarily $\textstyle 2^m < 2^n -1$, and thus the remainder in question is $\textstyle 2^m$. We want to show that $\textstyle 2^m \pmod {2^n-1}$ is an odd power of 2 for some $\textstyle n$, and thus not a power of 4.

 Let $\textstyle n=p^2$ for some odd prime $\textstyle p$. Then $\textstyle \varphi(p^2) = p^2 - p$. Since 2 is co-prime to $\textstyle p^2$, we have
 \[{2^{\varphi(p^2)} \equiv 1 \pmod{p^2}}\]and thus
 \[\textstyle 2^{p^2} \equiv 2^{(p^2 - p) + p} \equiv 2^p \pmod{p^2}.\]

 Therefore, for a counter-example, it suffices that $\textstyle 2^p \pmod{p^2}$ be odd. Choosing $\textstyle p=5$, we have $\textstyle 2^5 = 32 \equiv 7 \pmod{25}$. Therefore, $\textstyle 2^{25} \equiv 7 \pmod{25}$ and thus
 \[\textstyle 2^{2^{25}} \equiv 2^7 \pmod {2^{25} - 1}.\]Since $\textstyle 2^7$ is not a power of 4, we are done.

 



\textbf{Solution 2}

Lemma (useful for all situations): If $x$ and $y$ are positive integers such that $2^x - 1$ divides $2^y - 1$, then $x$ divides $y$.Proof: $2^y \equiv 1 \pmod{2^x - 1}$. Replacing the $1$ with a $2^x$ and dividing out the powers of two should create an easy induction proof which will be left to the reader as an Exercise.

 Consider $n = 25$. We will prove that this case is a counterexample via contradiction.

 Because $4 = 2^2$, we will assume there exists a positive integer $k$ such that $2^{2^n} - 2^{2k}$ divides $2^n - 1$ and $2^{2k} < 2^n - 1$. Dividing the powers of $2$ from LHS gives $2^{2^n - 2k} - 1$ divides $2^n - 1$. Hence, $2^n - 2k$ divides $n$. Because $n = 25$ is odd, $2^{24} - k$ divides $25$. Euler's theorem gives $2^{24} \equiv 2^4 \equiv 16 \pmod{25}$ and so $k \ge 16$. However, $2^{2k} \geq 2^{32} > 2^{25} - 1$, a contradiction. Thus, $n = 25$ is a valid counterexample.

 


\end{enumerate}